\chapter{Abstract}

% The abstract is a self-contained summary of the entire thesis (roughly 150--250 words).
% It must be understandable without reading the rest of the thesis.
% Do not use citations, unexplained abbreviations, or heavy mathematical notation.
%
% Follow the hourglass structure illustrated below: start broad (research area),
% narrow to the specific contribution, then broaden again to implications.

\begin{figure}[H]
\centering
\begin{tikzpicture}[
    every node/.style={font=\sffamily\small},
    label node/.style={font=\sffamily\footnotesize, text=white, align=center,
                       text width=3.8cm},
    arrow/.style={-{Stealth[length=5pt]}, thick, BrickRed!70!black}
]

%% --- hourglass geometry ---
% The hourglass is built from two filled trapezoids (top and bottom half)
% plus a thin waist rectangle, all in a matching color scheme.

% coordinates (x,y); hourglass is 6 cm wide at top/bottom, 1.4 cm at waist
\def\HW{3.0}   % half-width at top / bottom
\def\WW{1.3}   % half-width at waist
\def\TH{3.8}   % height of each trapezoid half
\def\RH{0.40}  % half-height of waist rectangle

% ---- top trapezoid (broad -> narrow) ----
\fill[BrickRed!20!white, draw=BrickRed!70!black, line width=0.6pt]
    (-\HW, \RH+\TH) -- (\HW, \RH+\TH)   % top edge
    -- (\WW, \RH)   -- (-\WW, \RH)       % bottom edge (waist)
    -- cycle;

% ---- bottom trapezoid (narrow -> broad) ----
\fill[BrickRed!20!white, draw=BrickRed!70!black, line width=0.6pt]
    (-\WW,-\RH) -- (\WW,-\RH)            % top edge (waist)
    -- (\HW,-\RH-\TH) -- (-\HW,-\RH-\TH) % bottom edge
    -- cycle;

% ---- waist rectangle ----
\fill[BrickRed!70!black]
    (-\WW,-\RH) rectangle (\WW,\RH);

%% --- labels inside the hourglass bands ---
% Top half: three bands (broad -> narrower -> waist)
\node[label node, text=BrickRed!80!black]
    at (0, \RH+\TH*0.82) {\textbf{Research Area}\\[1pt]
                           broad context,\\motivation};

\node[label node, text=BrickRed!80!black]
    at (0, \RH+\TH*0.50) {\textbf{Specific Problem}\\[1pt]
                           gap in the literature};

\node[label node, text=BrickRed!80!black]
    at (0, \RH+\TH*0.18) {\textbf{Approach}\\[1pt]
                           key idea};

% Waist: the contribution
\node[label node, text=white, font=\sffamily\footnotesize\bfseries,
      text width=2.4cm]
    at (0, 0) {Main Contribution};

% Bottom half: three bands (waist -> broader -> broad)

\node[label node, text=BrickRed!80!black]
    at (0, -\RH-\TH*0.30) {\textbf{Evaluation}\\[1pt]
                             implementation \& experiments};

\node[label node, text=BrickRed!80!black]
    at (0, -\RH-\TH*0.72) {\textbf{Implications}\\[1pt]
                             impact \& future directions};

%% --- down-arrows on the left side ---
\draw[arrow] (-\HW-0.5, \RH+\TH*0.60)
          -- (-\HW-0.5, \RH+\TH*0.28)
    node[midway, left, font=\sffamily\tiny, text=gray] {narrow};

\draw[arrow] (-\HW-0.5, \RH+0.1)
          -- (-\HW-0.5, -\RH-0.1)
    node[midway, left, font=\sffamily\tiny, text=gray] {focus};

\draw[arrow] (-\HW-0.5, -\RH-\TH*0.28)
          -- (-\HW-0.5, -\RH-\TH*0.60)
    node[midway, left, font=\sffamily\tiny, text=gray] {broaden};

%% --- brace labels on the right ---
\draw[decorate, BrickRed!60!black]
    (\HW+0.2, \RH+\TH) -- (\HW+0.2, \RH)
    node[midway, right=6pt, font=\sffamily\tiny, align=left,
         text=BrickRed!70!black] {motivation\\paragraph};

\draw[decorate, BrickRed!60!black]
    (\HW+0.2, -\RH) -- (\HW+0.2, -\RH-\TH)
    node[midway, right=6pt, font=\sffamily\tiny, align=left,
         text=BrickRed!70!black] {results \&\\discussion};

\end{tikzpicture}
\caption*{Hourglass structure of a good abstract.
          Start broad (research area and problem), funnel down to the specific
          contribution, then widen again to state results and implications.}
\end{figure}

\bigskip

Write a self-contained summary of roughly 150--250~words following the hourglass
structure above.
Cover the following five points in order:
\begin{enumerate}
    \item \textbf{Research area.}
          Name the field (e.g., automated termination analysis, probabilistic
          verification, complexity analysis of rewrite systems) and state in one
          sentence what the field studies and why it matters.
    \item \textbf{Specific problem.}
          State precisely which problem this thesis addresses and why it is
          challenging or not yet fully solved by existing approaches.
    \item \textbf{Approach.}
          Describe the key idea of your contribution in one or two sentences.
          Avoid mathematical notation and abbreviations, and use informal language.
    \item \textbf{Theoretical results.}
          Name the main property proved (soundness, completeness, an equivalence
          result, or a complexity bound) and state it informally.
    \item \textbf{Implementation and evaluation (if applicable).}
          Name the tool and the benchmark set (e.g., the Termination Problems
          Data Base) and give the headline result: how many more examples are
          solved compared to the baseline and what is the runtime overhead.
    \item \textbf{Implications.}
          Explain how your results solved the issue or fille the gap you pointed 
          out beforehand.
\end{enumerate}